\begin{figure}
  \centering
  \resizebox{\columnwidth}{!}{%
  \begin{tikzpicture}[
      font=\small\sffamily,
      tbl/.style={draw, rounded corners=2pt, align=center, minimum width=1.9cm,
                  minimum height=0.7cm, font=\footnotesize\sffamily},
      stp/.style={draw, rounded corners=2pt, align=center, fill=black!4,
                  minimum width=2.6cm, minimum height=1.1cm, font=\scriptsize\sffamily},
      cl/.style={draw, rounded corners=2pt, align=center, minimum width=1.9cm,
                 minimum height=0.8cm, font=\scriptsize\sffamily},
      jn/.style={draw, rounded corners=2pt, align=center, fill=black!8,
                 minimum width=1.9cm, minimum height=0.8cm, font=\scriptsize\sffamily},
      red/.style={draw, rounded corners=2pt, align=center, minimum width=5.4cm,
                  minimum height=0.9cm, font=\scriptsize\sffamily},
      outp/.style={draw, rounded corners=2pt, align=center, minimum width=1.9cm,
                   minimum height=0.6cm, font=\footnotesize\sffamily},
      flow/.style={-{Latex[length=2mm]}, thick},
      setup/.style={-{Latex[length=1.6mm]}, black!45},
      grp/.style={font=\scriptsize\bfseries\sffamily, rotate=90},
    ]

    % ---- inputs ----
    \node[tbl] (corpus) at (1.5,10.2) {corpus};
    \node[tbl] (probe)  at (4.5,10.2) {probe};

    % ---- steps 1 & 2 ----
    \node[stp] (s1) at (1.5,8.4) {\textbf{Step 1}\\ cluster corpus\\ k-means into $N$ clusters};
    \node[stp] (s2) at (4.5,8.4) {\textbf{Step 2}\\ cocluster probe\\ each row to $p$ nearest clusters};

    % ---- clusters ----
    \node[cl] (c1) at (0,6.4) {cluster 1\\ \scriptsize corpus + probe};
    \node[cl] (c2) at (3,6.4) {cluster 2\\ \scriptsize corpus + probe};
    \node[cl] (c3) at (6,6.4) {cluster 3\\ \scriptsize corpus + probe};

    % ---- step 3: exact join, per cluster ----
    \node[jn] (j1) at (0,4.5) {\textbf{Step 3}\\ (in-cluster) exact join};
    \node[jn] (j2) at (3,4.5) {\textbf{Step 3}\\ (in-cluster) exact join};
    \node[jn] (j3) at (6,4.5) {\textbf{Step 3}\\ (in-cluster) exact join};

    % ---- step 4: reduction (local + global) ----
    \node[red] (r) at (3,2.6)
      {\textbf{Step 4}\\ local/global reduction};

    % ---- output ----
    \node[outp] (out) at (3,1.1) {result};

    % ---- setup edges (clustering fan) ----
    \draw[setup] (corpus) -- (s1);
    \draw[setup] (probe)  -- (s2);
    \foreach \c in {c1,c2,c3} {
      \draw[setup] (s1) -- (\c);
      \draw[setup] (s2) -- (\c);
    }

    % ---- main flow ----
    \draw[flow] (c1) -- (j1);
    \draw[flow] (c2) -- (j2);
    \draw[flow] (c3) -- (j3);
    \draw[flow] (j1) -- (r);
    \draw[flow] (j2) -- (r);
    \draw[flow] (j3) -- (r);
    \draw[flow] (r)  -- (out);

    % ---- grouping boxes ----
    \node[draw, dashed, rounded corners, inner sep=8pt, fit=(j1)(j3)(r)] (exactbox) {};
    \node[draw, dotted, thick, rounded corners, inner sep=12pt,
          fit=(s1)(s2)(c1)(c3)(j1)(j3)(r)] (approxbox) {};

    % ---- group labels in the side margins (no arrow/box overlap) ----
    \node[grp] at ([xshift=-13pt]approxbox.west) {Approximate join (steps 1--4)};
    \node[grp] at ([xshift=13pt]exactbox.east)   {Exact join (steps 3--4)};
  \end{tikzpicture}}
  \caption{The vector join in four steps. Steps~3--4, the in-cluster
    brute-force join and its reduction, are the \emph{exact} join; the
    \emph{approximate} join wraps them in steps~1--2, which k-means cluster the
    corpus and route each probe row to its $p$ nearest clusters so the join runs
    per cluster. The exact join of Section~\ref{sec:technical-1:-exact} is this
    inner box over a single cluster, the whole table: every probe batch is
    brute-force joined against every corpus batch, then reduced. The reduction
    (step~4) runs pushed down inside each cluster (per-row top-$k$,
    max-distance) and again globally across clusters (global top-$k$).}
  \label{fig:exact}
\end{figure}
